\documentclass{article}
\usepackage[UTF8]{ctex}
\usepackage{amsmath,amssymb,amsfonts}
\usepackage{enumitem}
\usepackage{geometry}
\geometry{a4paper, margin=1in}
% === 错题本排版优化（不参与题目标识） ===
\setlength{\parindent}{0pt}
\setlength{\parskip}{0.35em plus 0.10em minus 0.08em}
\setlength{\emergencystretch}{3em}
\setlist{itemsep=0.25em, topsep=0.40em, parsep=0pt}
\allowdisplaybreaks[2]
% === 错题本排版优化结束 ===
\begin{document}
\section*{西安电子科技大学2024年数学分析真题}
\section*{190．西安电子科技大学2024年数学分析试卷}

\subsection{题目1}
1．求极限 \(\lim _{n \rightarrow \infty} \sqrt[n]{1+a^{n}+\left(\frac{a^{2}}{2}\right)^{n}}(a>0)\) ．

\subsection{题目2}
2．求不定积分 \(\int e^{-\sin x} \frac{\sin 2 x}{\sin ^{4}\left(\frac{\pi}{4}-\frac{x}{2}\right)} \mathrm{d} x\) ．

\subsection{题目3}
3．求极限 \(\lim _{x \rightarrow 0} \frac{6(\sqrt[3]{2-\cos x}-1)+\ln \left(1-x^{2}\right)}{x^{4}}\) ．

\subsection{题目4}
4．求重积分 \(\iint_{D} e^{\frac{y}{x+y}} \mathrm{~d} x \mathrm{~d} y\) ，其中 \(D\) 是 \(x+y=1, x=0, y=0\) 所围成的区域．

\subsection{题目5}
5．求曲线积分 \(\oint_{L} \frac{x \mathrm{~d} y-y \mathrm{~d} x}{4 x^{2}+y^{2}}\) ，其中 \(L\) 是圆周 \((x-1)^{2}+y^{2}=R^{2}\) ，取逆时针方向．

\subsection{题目6}
6．设 \(\sum\) 是 \(z=2-x^{2}-y^{2}\) 介于 \(z=1\) 和 \(z=2\) 之间的部分，取上侧，求曲面积分

\[
\iint_{\Sigma}\left(x^{3} z+x\right) \mathrm{d} y \mathrm{~d} z-x^{2} y z \mathrm{~d} z \mathrm{~d} x-x^{2} z^{2} \mathrm{~d} x \mathrm{~d} y .
\]

二．判断下列命题是否正确，若正确请给出证明，若错误则给出反例。每题6分，共24分。\\

\subsection{题目7}
1．若数列 \(\left\{a_{n}\right\}\) 以 0 为极限，数列 \(\left\{b_{n}\right\}\) 收敛，则级数 \(\sum_{n=1}^{\infty}\left(a_{n} b_{n}\right)\) 收敛。\\

\subsection{题目8}
2．设 \(f: \mathbb{R} \rightarrow \mathbb{R}\) 是连续映射，\(f_{n}=f \circ f \circ \cdots \circ f\) 为 \(f\) 的 \(n\) 次复合，若存在 \(x_{0} \in \mathbb{R}\) 使得数列 \(\left\{f_{n}\left(x_{0}\right)\right\}\)收敛到 \(a\) ，则 \(a\) 是 \(f\) 的不动点，即 \(f(a)=a\) 。\\

\subsection{题目9}
3．设 \(f(x), f_{n}(x)(n=1,2, \cdots)\) 均在 \([0,1]\) 上连续，且 \(\left\{f_{n}(x)\right\}\) 在 \([0,1]\) 上收敛于 \(f(x)\) ，则 \(\left\{f_{n}(x)\right\}\) 在\([0,1]\) 上一致收敛于 \(f(x)\) 。\\

\subsection{题目10}
4．设 \(D\) 是凸区域，若在 \(D\) 内满足 \(f_{x}(x, y)=0\) ，则 \(f(x, y)\) 在 \(D\) 内的取值与 \(x\) 无关．

三．（12 分）讨论 \(a, b, c\) 取何值时，\(\sum_{n=3}^{\infty} \frac{a^{n}}{n^{b}(\ln n)^{c}}\) 绝对收敛，条件收敛，发散．

四．（12 分）二元函数 \(f(x, y)\) 在点 \((1,0)\) 的某邻域有一阶连续偏导数，且 \(f(1,0)=0, f_{y}(1,0) \neq 0\) ，证明：方程 \(f\left(x, \int_{0}^{t} \cos u \mathrm{~d} u\right)=0\) 在点 \((1,0)\) 的某邻域内可确定隐函数 \(t=\varphi(x)\) ，并求 \(\varphi^{\prime}(1)\) 。

五。（15 分）设 \(f(x, y, z)=\ln x+2 \ln y+\ln z\) ，求 \(f(x, y, z)\) 在 \(x^{2}+y^{2}+z^{2}=4 r^{2}(x, y, z>0)\) 下的极大值，其中 \(r\) 为正常数。并证明对任意的正实数 \(a, b, c\) ，满足 \(a b^{2} c \leq 4\left(\frac{a+b+c}{4}\right)^{4}\) ．

六．（15 分）设 \(f(x)=\sum_{n=0}^{\infty} \frac{1}{2^{n}+x}\) ．

\begin{itemize}
  \item[（1）] 证明：\(f(x)\) 在 \([0,+\infty)\) 上可导且一致连续．
  \item[（2）] 证明：反常积分 \(\int_{0}^{+\infty} f(x) \mathrm{d} x\) 发散．
\end{itemize}

七．（15 分）设 \(f(x)\) 在 \([0,2]\) 上连续，\(h \in\left(0, \frac{1}{2}\right)\) ，记 \(f_{h}(x)=\frac{1}{h^{2}} \int_{0}^{h} \mathrm{~d} u_{2} \int_{0}^{h} f\left(x+u_{1}+u_{2}\right) \mathrm{d} u_{1}\) ．

\begin{itemize}
  \item[（1）] 证明：\(f_{h}(x)\) 在 \([0,1]\) 上二阶可导，并求 \(f_{h}^{\prime \prime}(x)\) ．
  \item[（2）] 若 \(f(x)\) 在 \(x_{0} \in[0,1]\) 处存在二阶导数，证明： \(\lim _{h \rightarrow 0^{+}} f_{h}^{\prime \prime}\left(x_{0}\right)=f^{\prime \prime}\left(x_{0}\right)\) ．
\end{itemize}

八．（15 分）设 \(f(x)\) 在 \([a, b]\) 上黎曼可积．

\begin{itemize}
  \item[（1）] 证明：对任意的 \(\varepsilon>0\) ，存在区间 \([c, d] \subset[a, b]\) ，使得 \(f(x)\) 在 \([c, d]\) 振幅 \(\omega_{f}<\varepsilon\) ．
  \item[（2）] 证明：\(f(x)\) 的连续点在 \([a, b]\) 上稠密，即对任意的 \([\alpha, \beta] \subset[a, b], f(x)\) 在 \([\alpha, \beta]\) 内部有连续点．
  \item[（3）] 若 \(f(x) \geq 0\) ，证明： \(\int_{a}^{b} f(x) \mathrm{d} x=0\) 当且仅当 \(f(x)\) 在其连续点恒为 0 ．
\end{itemize}
\end{document}
